\section{Discussion}

In this work we developed an algorithm for systematically computing and solving the loop equations
of general multi-matrix models, in the goal of finding the spectral curve for the planar resolvent
of a given matrix degree of freedom. We then focused our attention on the 3-state Potts model,
where we successfully computed the planar resolvents describing fixed, mixed, and free boundary
conditions using this technique. While the fixed and free boundary conditions were manifestly in
agreement with the literature, more care was required for the mixed boundary condition. Through
studying the analytic structure predicted by the saddle-point solution, we were able to show that
this too was consistent with the results of Atkin et al. We then considered some non-standard
resolvents, including a new `odd' boundary condition, which was susceptible to the loop equation
strategy, and a $U$ boundary condition which we were not able to solve. Following this, we
discussed the dilute Ising and Potts models, where we computed the resolvents corresponding to
various boundary conditions, in agreement with the results of Zinn-Justin.

There are many ways to continue this work. Our procedure is algorithmic in nature and so it would
be interesting to apply this to a wide variety of models for which the planar resolvent is known or
unknown. In the context of three-matrix models more generally, it would be interesting to
understand the criteria under which a resolvent is calculable. As we have seen from our results, we
were able to calculate the planar resolvent for particular systems with full $S_3$ symmetry under
interchange of the matrix degrees of freedom \eref{mixedAction}, and for systems where there is a
residual $\mathbb{Z}_2$ symmetry under the remaining degrees of freedom, with an auxiliary
condition reducing the number of level-1 resolvent functions \eref{freeAction}. Determining whether
this can be elevated to a general principle governing solvability of correlators in the
three-matrix model would be an exciting challenge. One immediate offshoot of this would be the
calculation of planar resolvents describing boundary conditions that interpolate between the fixed,
mixed and free boundary states analogous to \cite{Carroll:1997tr,Atkin:2012fx}. This would allow us
to directly study the boundary renormalisation group flow relating the different boundary states,
which is expected to induce a partial ordering of the spectrum of boundary states, conforming with
the boundary analogue of the $c$-theorem \cite{Zamolodchikov:1986gt}, which was first conjectured
in \cite{Affleck:1991tk}, and subsequently proven by Friedan and Konechny \cite{Friedan:2003yc}.

Another interesting avenue of research would be to investigate whether we could generalise the
algorithm developed in this chapter for a wider variety of correlators and matrix models, such as
the $q$-state Potts model where $q$ is taken to be non-integer. This is possible for the simple
example of the fixed-spin boundary condition \cite{Bonnet:1999nf}. Whether we could do the same for
general boundary conditions and reproduce the saddle-point results of Chapter 3 is an open
question. This would provide further insight into the nature of the allowed $G$-functions and
corresponding boundary conditions in the $q$-state Potts model. It would also be useful for
studying the expression for the New boundary condition determined in Chapter 4, which is given by a
non-standard resolvent function that would not immediately be amenable to calculation through the
loop equation method devised in this chapter.

Finally, it would be fruitful to extend our work to the calculation of observables in multi-matrix
models beyond the planar limit. It is well known that the planar resolvent forms the initial data
for higher-genus and multi-boundary correlators in the one- and two-matrix models
\cite{Eynard:2002kg,Eynard:2004mh,Eynard:2008we}, and that this may be extended to certain
multi-matrix models \cite{Borot:2009ia}. It would be interesting to see whether one could reframe
the method of loop equations for various boundary conditions in the 3-state Potts model such that
one can calculate higher-order correlation functions in an analogous topological recursion
procedure. This would allow, for example, the computation of cylinder amplitudes for different
boundary states, and might provide valuable information as to how to interpret the results of
Chapter 3.